\section{How we ran the experiments}
\label{sec:method}

\paragraph{Test set.} All comparisons use a fixed set of 23 incidents: one per fault
type per application (Train Ticket lacks one fault type). The wider corpus of 93
labelled incidents backs the larger baseline sweeps.

\paragraph{What we score.} Two things per diagnosis: did the agent name the right
component (\emph{localization}), and did it also pick the right fault type from the
menu (\emph{fully correct})? We report both throughout; localization alone is what
existing tools can be compared on.

\paragraph{What we compare against.} (1) A hand-written statistical method using the
same tools; (2) a published, peer-reviewed academic method (multi-source BARO from the
RCAEval project) run via its own released code; (3) two trace-only academic methods.
All see the same incident windows.

\paragraph{Degrading the data.} To ask ``how much observability is enough'', we
re-run the \emph{same} agent on deterministically thinned copies of the same
recordings: keeping 100/50/25/10/5\% of traces; sampling metrics every 5/10/30/60
seconds; keeping all logs, warnings-and-up, or errors-only; and giving the kernel
layer in full, in reduced forms, or not at all. The agent never changes within a
comparison --- only the data does.

\paragraph{Repeatability and noise.} LLM outputs vary between runs. By re-running
identical experiments we measured that any single condition's headline number carries
roughly a $\pm 9$-point band. We therefore repeat borderline incidents three times,
prefer per-fault-type flip evidence over small aggregate deltas, and never quote a
difference inside the band as a finding.

\paragraph{Cost.} A diagnosis costs about one cent of LLM usage (roughly 25k input
tokens, half to two-thirds served from the provider's cache, and $\sim$500 output
tokens); a full 23-incident experiment costs \$1--2. The entire experimental program
--- about 650 diagnoses --- cost under \$10.
